% HD-G26. GPT-6 Astra Pro /root/g34; inherited effort not numerically exposed.
% 21 November 2026. Actual correlated scalar-permutation groups.
\subsection{HD-G26: correlated Gaussian deleted-permutation sources of period rank nine are split}
\label{sec:correlated-permutation}

Let \(N\ge3\), \(K=\mathbb Q(i)\), and let
\[
 G\subset\mu_4\times S_N,\qquad
 \chi:G\longrightarrow\mu_4,\qquad
 V=\{(x_1,\ldots,x_N)\in K^N:\ \sum x_j=0\},
 \quad \rho(g)=\chi(g)P_g|_V.                         \tag{G26.1}
\]
Assume this actual representation is absolutely irreducible and
has character field exactly \(K\). For a connected full-\(G\)
sphere cover, suppose its natural primitive factor is a Weil
sixfold of signature \((3,3)\), with the restricted theta
polarization and the Gaussian action constructed below.
Let \(r_0\) count the true branch values whose scalar character
is nontrivial. Its ordinary Hermitian determinant is
\[
                         [-N^{\,r_0-2}].                    \tag{G26.2}
\]
Odd \(r_0\) forces period rank at most seven. In particular
every such period-rank-nine family is split. No independence of
the two projections of \(G\) is assumed, and \(G\) need not
contain the scalar \(iI\). This includes the reflection-free
correlated groups in arbitrary rank \(N-1\).

\paragraph{1. The actual point stabilizer has the full scalar character.}
Write \(P\) for the permutation image of \(G\). Untwisting by
the one-dimensional character \(\chi\) shows that its ordinary
augmentation module is absolutely irreducible. It has dimension
at least two, so has no invariant vector. The dimension of
invariants in the permutation module is the number of \(P\)-orbits;
hence \(P\) is transitive. The permutation module is the direct
sum of its constant line and its absolutely irreducible augmentation
module. Its commuting endomorphism algebra consequently has
dimension two. A matrix commutes with the permutation action
precisely when its entries are constant on the orbits in the
ordered pair set. There are thus exactly two such orbits:
the diagonal and its complement. This proves two-transitivity
and, in particular, primitivity of the actual permutation action.
For the latter implication, a block containing two distinct points
is preserved by the stabilizer of either one; transitivity of
that stabilizer on the other points makes the block the whole set.

The image of \(\chi\) is all \(\mu_4\). Otherwise every trace
\(\chi(g)(\#\operatorname{Fix}(P_g)-1)\) would be rational,
contrary to the specified character field. Let
\[
 H=\operatorname{Stab}_G(1),\qquad
 K_0=\ker\chi,\qquad H_0=H\cap K_0.
                                                               \tag{G26.3}
\]
If \(\chi(H)\ne\mu_4\), then \(H\subset\ker\chi^2\), a
proper subgroup of index two in \(G\). Intermediate subgroups
between \(H\) and \(G\) give blocks in the actual action on
\(G/H\): the block is the set of cosets contained in that
subgroup. Primitivity therefore forces \(H=\ker\chi^2\),
giving \(N=2\), a contradiction. Consequently
\[
 \chi(H)=\mu_4,\qquad [H:H_0]=4,\qquad
 |H|=|G|/N,\qquad |K_0|=N|H_0|.                         \tag{G26.4}
\]
Choose an actual \(h_0\in H\) with \(\chi(h_0)=i\).
It need not be central. Since \(HK_0=G\), the quotient
maps used below have connected fibre product.

The coefficient representation in (G26.1) is faithful. If a
scalar permutation acts trivially on the augmentation space,
its real permutation action there is a real scalar \(1\) or
\(-1\). Scalar \(-1\) would make its permutation trace
\(1-(N-1)=2-N<0\), impossible for a permutation matrix.
Scalar \(1\) gives trace \(N\), so the permutation is the
identity, and its scalar phase is then also one. Thus every
true branch element has positive coefficient codimension.

\paragraph{2. Exact coefficient splitting and integral boundary maps.}
Put \(M=K^N\) with action \(\chi(g)P_g\), and let
\(K_\chi\) have scalar action \(\chi\). The actual maps
\[
 j(a)=(a,\ldots,a),\qquad
 p(x)=N^{-1}\sum_jx_j,\qquad pj=1
                                                               \tag{G26.5}
\]
split \(M=K_\chi\perp V\) for \(h(x,y)=\sum\bar x_jy_j\).
The constant-line pullback is \(Nh_\chi\). In the basis
\(e_j-e_N\) of \(V\), the coefficient Gram matrix is
\(I+\boldsymbol1\boldsymbol1^t\), of determinant \(N\).

Use the self-dual coefficient lattices
\[
                  M_R=R^N,\quad R_\chi=R,\qquad
                         R=\mathbb Z[i,1/2].                \tag{G26.6}
\]
Global zeroth homology vanishes for both and for their duals.
Indeed \(h_0\) fixes \(e_1\) with multiplier \(i\); the
coinvariant relation \((i-1)e_1=0\) kills \(e_1\), since
\(i-1\) is a unit of \(R\). Transitivity then kills every
coordinate. The same argument uses \(-i-1\) for the dual.
For the scalar system, surjectivity of \(\chi\) gives the
same assertion. This is the replacement for a central
scalar in the full-product case.

For a local element \(g\), put \(\zeta=\chi(g)\).
On a permutation cycle of length \(\ell\), the unit basis
\(w_j=\zeta^je_{P_g^j(a)}\) gives
\(gw_j=w_{j+1}\) until the last step and
\(gw_{\ell-1}=\zeta^\ell w_0\). Its coinvariant module is
exactly \(R/(\zeta^\ell-1)\). If \(\zeta^\ell=1\),
it is free on \([w_0]\), and the invariant generator
\(\sum_jw_j\) pairs with it by one. The relations
\(w_{j+1}-w_j\) give the kernel of the coefficient-sum map.
If \(\zeta^\ell\ne1\), then \(\zeta^\ell-1\) is a unit
of \(R\), both modules are zero, and the literal inverse
to \(1-g\) is
\[
                     \frac{1+g+\cdots+g^{\ell-1}}
                                  {1-\zeta^\ell}.
                                                               \tag{G26.7}
\]
Thus all boundary invariant and coinvariant modules are free,
without division by cycle lengths, \(N\), or \(|G|\).
The scalar line is the length-one instance.

\paragraph{3. Perfect parabolic lattices and the determinant.}
Remove small discs about the branch points, giving a compact
surface with boundary \(\Sigma\). Absolute first homology of
either coefficient lattice in (G26.6) is free, from the graph
cellular model and the just-proved vanishing of absolute zeroth
homology. Integral Poincare--Lefschetz duality and the universal
coefficient sequence give perfect cross-pairings with relative
first homology, which is also free. The pair exact sequence
identifies the cokernel of the absolute-to-relative map with
the free boundary zeroth homology. Its image is therefore
primitive. The two such maps are adjoint. Extending a functional
from the primitive image to the relative module and applying
the perfect cross-pairing proves that the induced parabolic
pairing is perfect over \(R\). These are the actual cap-product
and restriction maps; no odd-order averaging has been used.
They identify the rational parabolic image with the middle
extension cohomology of the finite local system.
This is the integral primitive-image argument of HD-S27 and
HD-S33; the general cap-product comparison is
Dettweiler--Wewers, Proposition2.1, printed pp.153--154.

Let \(H_M=iS_M\), \(H_\chi=iS_\chi\) denote the associated
Hermitian parabolic forms in the inverse-deck convention.
Their determinants are rational units of \(R\), hence signed
powers of two. The scalar parabolic rank is \(r_0-2\),
by the sphere Euler formula and its two zero invariant spaces.
The assumed sixfold factor gives parabolic rank six for \(V\).
Applying (G26.5) to the actual cellular and absolute-to-relative
maps gives the orthogonal decomposition
\[
 (P(M)_K,H_M)\simeq
          (P(\chi)_K,NH_\chi)\perp(P(V),H_V),
 \qquad \operatorname{rank}_KP(M)=r_0+4.                 \tag{G26.8}
\]
In bases, its change-of-basis determinant contributes exactly a
\(K/\mathbb Q\) norm. Therefore
\[
 [\det H_V]=[\det H_M]\,
            [N^{r_0-2}\det H_\chi]^{-1}.                \tag{G26.9}
\]
The quotient of the two perfect-lattice determinants is a signed
power of two; the signature \((3,3)\) forces its sign to be
negative. Powers of two are Gaussian norms, and replacing
\(N^{-(r_0-2)}\) by \(N^{r_0-2}\) also changes by a norm.
This proves (G26.2) for the coefficient form. The next paragraph
identifies it with the actual restricted polarization.

\paragraph{4. The quotient fibre product and its actual complementary factor.}
Let \(C\to\mathbb P^1\) be the full cover and form the
actual quotient curves
\[
 Z=C/K_0,\qquad Y=C/H,\qquad X=C/H_0,\qquad
 \alpha:C\to X,\qquad f:X\to Z.                          \tag{G26.10}
\]
Their degrees over the sphere are \(4,N,4N\). Label a
regular cover fibre by the left deck action, so quotient
fibres are left coset sets. The map
from \(H_0\backslash G\) to
\((H\backslash G)\times(K_0\backslash G)\) is injective
and is surjective because both sets have size \(4N\);
equivalently use \(HK_0=G\) and \(H\cap K_0=H_0\).
Over the branch complement it is the actual map of cover
fibres. Consequently \(X\) is the normalization of
\(Y\times_{\mathbb P^1}Z\). This verifies the fibre product
also at its ramified points.

The element \(h_0\) normalizes \(H_0\), so it descends
to the order-four deck transformation \(\sigma\) of \(X/Y\)
and to its phase-\(i\) deck transformation on \(Z\).
The degree-\(N\) map \(f\) commutes with these two cyclic
actions. Pullback and norm give \(f_*f^*=[N]\) and are
adjoint for the principal curve Jacobian polarizations.
Put
\[
 \epsilon_X=(1-a_\sigma^2)/2,\qquad
 \pi_X=\epsilon_X(1-N^{-1}f^*f_*),\qquad
 A_X=\operatorname{im}(2N\pi_X:JX\to JX).
                                                               \tag{G26.11}
\]
This is a self-adjoint rational idempotent; its displayed
multiple is integral. Here is the literal local coefficient
map, including its action convention. Realize the induced
module by functions \(F:G\to K\) with
\(F(hx)=\chi(h)F(x)\) for \(h\in H\), with right
translation \((gF)(x)=F(xg)\). Then
\[
 v\longmapsto F_v(x)=(\rho_M(x)v)_1
             =\chi(x)v_{P_x^{-1}(1)}
                                                               \tag{G26.11a}
\]
is an equivariant isomorphism from \(M\): its values on
coset representatives recover every coordinate, with only
unit phases. On \(H_0\backslash G\) it is exactly the
odd cyclic-sheet function of \(X\). Together with its
conjugate it identifies the primitive cyclic coefficient
system of \(X\) with \(M\) and \(\bar M\).
The pullback from \(Z\) is the constant-sheet
map \(j\) in (G26.5); its trace is the sum of the \(N\)
sheets. Their composite is exactly \(N\). Filling each
branch point uses the invariant stalks computed in paragraph2.
It follows that (G26.11) has exactly the paired augmentation
multiplicity space and hence dimension six. Its Gaussian
action is the actual inverse cyclic action \(a_\sigma^{-1}\).

\paragraph{5. A literal primitive idempotent and the exact isogeny.}
Take \(w=Ne_1-\sum_je_j\), of norm
\(q=N(N-1)\). It satisfies \(\rho(h_0)w=iw\).
Define
\[
 e=\frac{2}{|G|N}\sum_{g\in G}
       \operatorname{Re}h(w,\rho(g^{-1})w)\,g,\qquad
 D=\frac{|G|N}{2}.                                      \tag{G26.12}
\]
The integer \(D\) exists because \(4\mid |G|\).
The coefficients of \(De\) are integers. Schur
orthogonality, absolute irreducibility and the distinct
conjugate character identify \(e\) with the rank-one
orthogonal projector onto \(Kw\) in \(V\), its conjugate
projector in \(\bar V\), and zero on every other
irreducible. Thus it is a primitive rational self-adjoint
idempotent with corner \(K\).

Here the entire comparison can also be made coefficientwise.
The exact inner product is
\[
 h(w,\rho(g^{-1})w)
       =\chi(g)^{-1}N\bigl(N\boldsymbol1_H(g)-1\bigr).
                                                               \tag{G26.13}
\]
For a subgroup \(L\), let
\(\operatorname{av}_L=|L|^{-1}\sum_{\ell\in L}\ell\).
Writing the four cosets of \(H_0\) in \(H\), and of
\(K_0\) in \(G\), gives
\[
 \begin{split}
 e&=e_{H,\chi+\bar\chi}-e_{G,\chi+\bar\chi}\\
  &=\frac{1-h_0^2}{2}
                \bigl(\operatorname{av}_{H_0}
                               -\operatorname{av}_{K_0}\bigr).
 \end{split}                                                   \tag{G26.14}
\]
Indeed the coefficient of \(g\) in the first line is
\(2|G|^{-1}(N\boldsymbol1_H(g)-1)
\operatorname{Re}\chi(g)^{-1}\), exactly (G26.12)--(G26.13).
Although \(h_0\) need not be central in \(G\), it
commutes with both averages in (G26.14). Also \(eh_0^2=-e\).

Use left deck maps \(L_g\), the covariant Jacobian maps
\(a_g=L_{g,*}\), and put \(B=\operatorname{im}(a_{De})\).
Then \(a_{h_0}\) preserves \(B\), and
\[
                          i_B=a_{h_0}^{-1}|_B             \tag{G26.15}
\]
is an actual integral Gaussian action. Its Rosati adjoint
is its negative, since deck inverses are Rosati adjoints
and its square is \(-1\).

The quotient identities
\(\alpha^*\alpha_*=\sum_{h\in H_0}a_h\) and
\((f\alpha)^*(f\alpha)_*=\sum_{k\in K_0}a_k\), together
with \(|K_0|=N|H_0|\), turn (G26.14) into the literal
equality of rational Jacobian projectors
\[
                \frac1{|H_0|}\alpha^*\pi_X\alpha_*=a_e.
                                                               \tag{G26.16}
\]
Thus \(\alpha^*(A_X)\) has exactly the rational homology
image of \(a_e\); the connected abelian subvarieties coincide.
The actual homomorphism
\[
                  q_X=\alpha^*|_{A_X}:A_X\longrightarrow B
                                                               \tag{G26.17}
\]
is an isogeny. It intertwines (G26.15) and the inverse cyclic
action on \(A_X\), because \(\alpha\) intertwines \(h_0\)
and \(\sigma\). Its polarization scalar is exactly
\[
             q_X^*(\theta_C|_B)
                        =|H_0|\,(\theta_X|_{A_X}).        \tag{G26.18}
\]
This is the usual pullback/norm adjunction for the actual
degree-\(|H_0|\) map \(\alpha\), now restricted to the
identified factors. No unspecified projection scalar enters.

To fix the coefficient normalization as well, put
\[
 j_x(g)=h(w,\rho(g)x)/q,\qquad
                      \Phi(x)=\operatorname{PD}(j_x+\bar j_x).
 \]
For the cohomology action \(T_g=(L_g^{-1})^*\), one has
\(T_{h_0}j_x=j_{iw,x}=-ij_x\). Hence
\(a_{h_0}\Phi(x)=\Phi(-ix)\), which proves exactly the
\(K\)-linearity of \(\Phi\) for (G26.15).
The complex Schur identity and the two conjugate cross terms
give the exact Hermitian multiplier
\[
 h_B(\Phi x,\Phi y)
          =c\,H_V(x,y),\qquad
 c=\frac{2|G|}{(N-1)q}
                    =\frac{2|G|}{N(N-1)^2}.               \tag{G26.19}
\]
All these are the matrix-coefficient and Poincare-duality
maps of HD-S23, with the vector norm explicitly retained.
Clearing their homology denominators gives the corresponding
actual isogenies. The positive rational multiplier changes
the rank-six determinant by \(c^6=N_{K/\mathbb Q}(c^3)\).
Equations(G26.18)--(G26.19) prove (G26.2) for both actual
restricted polarizations.

\paragraph{6. Odd scalar count and the actual period-rank bound.}
For a local permutation \(P_g\), let \(k\) be its cycle
count and \(k_a\) the number of cycles whose lengths are
divisible by \(a\). The fixed-vector equations on each
cycle, with the scalar constant line removed, give
\[
 c_V(g)=
 \begin{cases}
 N-k,&\chi(g)=1,\\
 N-1-k_2,&\chi(g)=-1,\\
 N-1-k_4,&\chi(g)=i\text{ or }-i.
 \end{cases}                                               \tag{G26.20}
\]
The generating branch phases generate \(\mu_4\), and their
product is one. Thus the number of odd phases is even and
at least two. If \(r_0\) is odd, at least one additional
phase is \(-1\). Select two odd phases and one \(-1\).
Their total codimension is at least
\[
 2(N-1-\lfloor N/4\rfloor)
       +N-1-\lfloor N/2\rfloor\ge2N-3,
 \]
since \(2\lfloor N/4\rfloor+\lfloor N/2\rfloor\le N\).
Every other true branch contributes at least one by faithfulness.
The parabolic sixfold equality is
\(\sum c_V=2(N-1)+6=2N+4\). If \(b\) is the number
of true branch values, it follows that
\[
                    2N+4\ge(2N-3)+(b-3),\qquad b\le10.
                                                               \tag{G26.21}
\]
For fixed finite data, the marked cover family is locally
finite over its ordered branch configuration: it has finitely
many generating tuples, and continuation on a cut surface
followed by the fixed local branched-disc maps gives its
actual deformation. After quotienting by \(PGL_2\), its
dimension is \(b-3\). The relative Jacobian and the fixed
projector (G26.12) give its period morphism, whose differential
rank is therefore at most seven. No independence of the
permutation and scalar projections is used.

Even \(r_0\) makes (G26.2) the ordinary class \([-1]\).
The explicit Meyer reductions and final binary norm equation
of HD-S03a identify a signature-\((3,3)\) space with this
class with \(U^3\). This is the asserted split conclusion.

\paragraph{Scope and finite checks.}
The proof concerns the specified actual deleted-permutation
representation and its actual restricted-theta primitive factor.
It is stable under a specified \(K\)-similitude or the
fixed-isogeny flat-polarization comparisons of HD-S17 on a
rank-nine family. It makes no assertion about arbitrary
representations of the same abstract group or arbitrary
polarizations at a single special point.

The companion finite computation uses the actual correlated
group obtained by the quartic determinant character on
\(\operatorname{AGL}_2(\mathbb F_5)\), acting on its
twenty-five points. It rebuilds the group, its full stabilizer
character and noncentral \(h_0\), the literal projector
identity, and the coefficient Schur normalization. This is
a finite supplement to the universal proof. No sphere cover
with the required sixfold multiplicity is asserted for that
example.
